\subsection{Discrete PID with filtered derivative}
The controller is a parallel PID sampled at $\Delta t = 0.02$~min, with $u_{\text{raw},k} = u_{\text{bias}} + P_k + I_k + D_k$, $u_{\text{bias}} = q_{j,ss}$ and $e_k = y_{sp,k} - y_{pv,k}$ \cite{astrom2006pid, seborg2010process}:
\begin{align}
    P_k &= K_p e_k, \label{eq:proportional_term} \\
    I_k &= I_{k-1} + \frac{K_p}{T_i}\,\frac{e_k + e_{k-1}}{2}\,\Delta t, \label{eq:integral_trapezoidal} \\
    D_k &= \frac{\tau_f}{\tau_f + \Delta t} D_{k-1} - \frac{K_p T_d}{\tau_f + \Delta t}\,(y_{pv,k} - y_{pv,k-1}), \label{eq:derivative_filtered}
\end{align}
with derivative on the measurement and filter time $\tau_f = T_d/N$ \cite{astrom1995pid, visioli2006practical}. All loops use the common value $N = 10$; Section~\ref{subsec:measurement_path} discusses its effect on the noise-to-valve gain \cite{segovia2013noise}. The loop is reverse-acting: more coolant lowers the temperature, so all tunings have $K_p < 0$.

\subsection{Directional amplitude- and rate-aware clamping}
\label{subsec:clamping_rule}
Let $u_{\text{act},k-1}$ be the valve position actually reached at the previous sample after amplitude \emph{and} slew-rate limitation. With the trapezoidal increment $\delta I_k = (K_p/T_i)\,\tfrac{1}{2}(e_k + e_{k-1})\,\Delta t$ of Eq.~\eqref{eq:integral_trapezoidal}, the integrator is frozen when the previous command was cut by either limit and the increment pushes further in the same direction:
\begin{equation}
    I_k = \begin{cases}
        I_{k-1}, & \text{if } \left(u_{\text{raw},k-1} - u_{\text{act},k-1}\right)\delta I_k > 0, \\
        I_{k-1} + \delta I_k, & \text{otherwise.}
    \end{cases}
    \label{eq:clamping_antiwindup}
\end{equation}
Evaluating the saturation residual at the previous sample is an implementation choice, not a necessity: it needs no iteration, at the price of a one-sample delay of the freeze decision (the integrator can still move by one increment $\delta I_k$ after a limit has cut the command). A predictor-corrector test at the current sample (Supplementary Section~\suppref{supp:aw}, Table~\suppref{tab:predictive_clamping}) removes the delay; on the nominal step of Table~\ref{tab:antiwindup_ablation} it changes the IAE by at most $\vpcnommaxpct\%$, and on all steps by at most $\vpcrestmaxpct\%$ for TL, AMIGO, SIMC and Optimal ITAE but by up to $\vpcZNmaxpct\%$ and $\vpcCCmaxpct\%$ for ZN and CC on the $+5$~K step, so for high-gain loops the delay is not negligible, although no conclusion of Section~\ref{sec:results} changes. The $k-1$ test remains the default because it needs only stored signals, no evaluation of the valve map inside the controller, and is the rule used for every result. The 'no extra parameter' property is relative to $S_{\max}$ and $\Delta t$, which enter through $u_{\text{act}}$ (Section~\ref{subsec:ablation}). The product form contains the sign of the loop: with $K_p < 0$ a negative error drives the command towards $u_{\max}$, so the sign-blind condition ``$u_{\text{raw}} \ge u_{\max}$ and $e_k > 0$'' is false while the integrator winds up and becomes true only when the error has reversed, i.e. when the integrator must unwind (Section~\ref{subsec:ablation}). It also covers rate saturation: a $2$~K setpoint step demands more than $S_{\max}\Delta t = 1.6$~L/min per sample from most tunings, so the valve spends the first part of the transient rate-limited but far from its amplitude stops, and an amplitude-only rule keeps integrating during the whole ramp \cite{bohn1995analysis, peng1996nonlinear}. Section~\ref{subsec:ablation} quantifies the difference against sign-blind clamping and the other variants defined next. The rule is the sign-of-increment form of conditional integration, evaluated on the slew-limited valve position; it is adopted, not proposed as new (Section~\ref{subsec:positioning}).


\begin{proposition}[latch of sign-blind clamping]\label{prop:latch}
Let the loop of Eqs.~\eqref{eq:proportional_term}--\eqref{eq:derivative_filtered} be reverse-acting ($K_p < 0$), with $u_{\text{act},k} = \operatorname{rate}(\operatorname{sat}(u_{\text{raw},k}))$ held over $[t_k, t_{k+1})$, and let the plant be open-loop stable with steady temperature $T_\infty = T_{eq}(u_{\max})$ at full cooling, $T_{sp} > T_\infty$. Sign-blind amplitude clamping sets $I_k = I_{k-1}$ if $u_{\text{raw},k-1} \ge u_{\max} \wedge e_k > 0$ or $u_{\text{raw},k-1} \le u_{\min} \wedge e_k < 0$\footnote{The released code also requires that a limit cut the previous command, $u_{\text{act},k-1} \ne u_{\text{raw},k-1}$; the two rules differ only when $u_{\text{raw},k-1}$ equals a stop exactly.}, and $I_k = I_{k-1} + \delta I_k$ otherwise. Let $I^\star = u_{\max} - u_{\text{bias}} + |K_p|\,(T_{sp} - T_\infty)$, $\delta > 0$, $\bar D \ge 0$ and assume
(H1) $u_{\text{act},k_0} = u_{\max}$ and $I_{k_0} \ge I^\star + |K_p|\delta + \bar D$ at some sample $k_0$;
(H2) the response to the constant input $u_{\max}$ from the state at $t_{k_0}$ has $T_k \in [T_\infty - \delta, T_{sp})$ for all $k > k_0$;
(H3) $|D_k| \le \bar D$ for all $k > k_0$.
Then $I_k = I_{k_0}$, $u_{\text{raw},k} \ge u_{\max}$ and $u_{\text{act},k} = u_{\max}$ for all $k \ge k_0$, and $T_k \to T_\infty$: the loop stays at the offset $e_\infty = T_{sp} - T_\infty > 0$.
\end{proposition}
\begin{proof}
By induction on $k \ge k_0$. Base: the slew limiter moves $u_{\text{act}}$ towards $\operatorname{sat}(u_{\text{raw},k_0})$ without passing it, so $u_{\text{act},k_0} = u_{\max}$ implies $u_{\text{raw},k_0} \ge u_{\max}$. Step: if the claim holds up to $k$, the input has been $u_{\max}$ on $[t_{k_0}, t_{k+1})$, so by (H2) $0 < e_{k+1} \le T_{sp} - T_\infty + \delta$. With $u_{\text{raw},k} \ge u_{\max}$ and $e_{k+1} > 0$ the rule freezes $I_{k+1} = I_k = I_{k_0}$, and by (H1) and (H3)
\begin{align*}
u_{\text{raw},k+1} &= u_{\text{bias}} - |K_p|\,e_{k+1} + I_{k_0} + D_{k+1} \\
&\ge u_{\text{bias}} - |K_p|(T_{sp} - T_\infty + \delta) \\
&\quad + I^\star + |K_p|\delta = u_{\max},
\end{align*}
so $\operatorname{sat}(u_{\text{raw},k+1}) = u_{\max} = u_{\text{act},k}$ and $u_{\text{act},k+1} = u_{\max}$. The input is constant from $t_{k_0}$ on and the plant is stable, hence $T_k \to T_\infty$.
\renewcommand{\qedsymbol}{}
\end{proof}
(H3) is a hypothesis, but it follows from (H2) and Eq.~\eqref{eq:derivative_filtered} (Supplementary Eq.~\suppref{eq:dbar}). (H1) and (H2), in contrast, are properties of a trajectory and are verified by simulation, not proved. On the reactor $I^\star = \vZNIstar$~L/min for ZN and $\vCCIstar$~L/min for CC; from $t_{k_0} = \vZNtkzero$ and $\vCCtkzero$~min the valve stays at $u_{\max}$, the band holds with $\delta = \vZNdelta$ and $\vCCdelta$~K, and $I_{k_0} = \vZNIkzero$ and $\vCCIkzero$~L/min exceed $I^\star + |K_p|\delta + \bar D$ even with the band bound $\bar D = \vZNDbar$ and $\vCCDbar$~L/min (\texttt{results/jpc/latch\_check.csv}); the latch is a property of the rule, not of the tuning. The directional rule of Eq.~\eqref{eq:clamping_antiwindup} cannot latch: at $u_{\text{act}} = u_{\max}$, $u_{\text{raw}} > u_{\max}$ and $e > 0$ the increment is $\delta I = (K_p/T_i)\tfrac12(e_k + e_{k-1})\Delta t < 0$, so $(u_{\text{raw}} - u_{\text{act}})\,\delta I < 0$, the integrator is \emph{not} frozen and unwinds by $|\delta I|$ per sample until $u_{\text{raw}} \le u_{\max}$. A fixed point with the valve at a stop requires $(u_{\text{raw}} - u_{\text{act}})\,\delta I > 0$, i.e.\ at the upper stop $e < 0$ (temperature above the setpoint, more cooling genuinely needed); every other fixed point has $e = 0$.

\subsection{Anti-windup variants of the ablation}
\label{subsec:aw_variants}
Seven integrator policies are compared (Table~\ref{tab:antiwindup_ablation}). \emph{None}: $I_k = I_{k-1} + \delta I_k$ at all times. \emph{Back-calculation} \cite{astrom1989windup, galeani2009tutorial} adds a tracking term,
\begin{equation}
    I_k = I_{k-1} + \delta I_k + \frac{\Delta t}{T_t}\left(u_{\text{act},k-1} - u_{\text{raw},k-1}\right),
    \label{eq:back_calculation}
\end{equation}
with the default $T_t = \sqrt{T_i T_d}$ ($T_t = T_i$ if $T_d = 0$); because the scheme depends on $T_t$, Section~\ref{subsec:backcalc} scans it over the physical range $[T_d, T_i]$ and beyond. \emph{Sign-blind amplitude clamping} freezes the integrator when $u_{\text{raw}} \ge u_{\max}$ and $e_k > 0$, or $u_{\text{raw}} \le u_{\min}$ and $e_k < 0$, which is the sign convention of a direct-acting loop. The remaining four variants form the $2\times 2$ factorial of the clamping family. The \emph{trigger} is either amplitude only ($u_{\text{raw},k-1}$ outside $[u_{\min}, u_{\max}]$) or amplitude plus slew rate ($u_{\text{raw},k-1} \ne u_{\text{act},k-1}$). The \emph{sign awareness} is either unconditional (the integrator is frozen whenever the trigger holds, whatever the sign of $\delta I_k$) or directional (frozen only if $(u_{\text{raw},k-1} - u_{\text{act},k-1})\,\delta I_k > 0$, Eq.~\eqref{eq:clamping_antiwindup}). The rule adopted in the benchmark is the amplitude-plus-rate, directional cell; Algorithm~\ref{alg:integrator} gives its implementation.

\begin{algorithm}[t]
\caption{Integrator update with directional amplitude- and rate-aware clamping (one sampling instant, reverse- or direct-acting).}\label{alg:integrator}
\begin{algorithmic}[1]
\Require $e_k$, $e_{k-1}$, $I_{k-1}$, previous command $u_{\text{raw},k-1}$ and previous valve position $u_{\text{act},k-1}$ (after amplitude \emph{and} slew limits)
\State $\delta I \gets (K_p/T_i)\,\tfrac12 (e_k + e_{k-1})\,\Delta t$
\If{$(u_{\text{raw},k-1} - u_{\text{act},k-1})\,\delta I > 0$}
    \State $I_k \gets I_{k-1}$ \Comment{limit cut the last move and the increment pushes further into it}
\Else
    \State $I_k \gets I_{k-1} + \delta I$
\EndIf
\State $u_{\text{raw},k} \gets u_{\text{bias}} + P_k + I_k + D_k$; \quad $u_{\text{act},k} \gets \operatorname{rate}(\operatorname{sat}(u_{\text{raw},k}))$
\end{algorithmic}
\end{algorithm}

\subsection{Checklist for reporting the anti-windup logic}
\label{subsec:checklist}
A comparison of controllers on a valve-limited loop is only interpretable if it states the following five items; each one changed a conclusion in this study (the order of evaluation in Section~\ref{subsec:robustness_ext}).
\begin{enumerate}
    \item \textbf{Loop sign.} Is the loop reverse-acting ($K_p<0$ in the convention used), and was the clamping condition written for that sign?
    \item \textbf{Trigger and order.} Does the logic react to the amplitude limit only, or also to the slew-rate limit (i.e.\ does it use the post-rate-limit valve position), and in which order are amplitude, slew and backlash evaluated?
    \item \textbf{Sign awareness.} Is the integrator frozen unconditionally, or only if the increment pushes further into the limit?
    \item \textbf{Tracking constant.} If back-calculation is used, what is $T_t$, is it inside $[T_d, T_i]$ and larger than $\Delta t$, and how sensitive is the result to it?
    \item \textbf{Measurement path.} Which derivative-filter coefficient $N$ and which measurement filter $\tau_m$ were used, how was the filter discretized, and are the errors evaluated on the true or on the measured signal?
\end{enumerate}

\subsection{Measurement filter}
\label{subsec:meas_filter}
Where indicated, the controller sees the reactor temperature through a first-order low-pass filter, $\tau_m\,\dot T_{\text{meas}} = T + \eta - T_{\text{meas}}$ with $\tau_m = 0.05$~min and sensor noise $\eta$ (the symbol $\tau_f$ is reserved for the derivative filter). Under the zero-order hold at the sampling instants it is integrated exactly, $T_{\text{meas},k} = T_{\text{meas},k-1} + (1 - e^{-\Delta t/\tau_m})\,(T_k + \eta_k - T_{\text{meas},k-1})$. Error metrics keep using the true temperature $T$. The filter adds lag inside the loop; the benchmark tables use $\tau_m = 0$ unless stated.

\subsection{FOPTD identification}
\label{subsec:foptd_identification}
Since $\mathbf{x}_{ss}$ is open-loop stable, a $+10$~L/min coolant step was applied to the nonlinear model and a first-order-plus-dead-time (FOPTD) model
\begin{equation}
    G_p(s) = \frac{K}{\tau s + 1}\,e^{-\theta s}
    \label{eq:foptd}
\end{equation}
was fitted by the two-point method of Sundaresan and Krishnaswamy followed by nonlinear least squares, giving $K = -0.1852$~K/(L/min), $\tau = 0.6427$~min, $\theta = 0.0935$~min (RMSE $0.0092$~K) and a lag-dominant ratio $r = \theta/\tau = 0.1455$ \cite{shinskey1996process, odwyer2009handbook}. The apparent delay lumps the two fast poles and the zero of $G_{\text{true}}$ in Eq.~\eqref{eq:g_true_exact}; its consequence for the sensitivity peak is shown by the two $M_s$ columns of Table~\ref{tab:testbed}.

\subsection{Test-bed gains}
\label{subsec:tuning_rules}
Six fixed-gain controllers span the range of aggressiveness for a lag-dominant loop; their gains are inputs to the anti-windup study and are not designed or compared as tuning methods:
\begin{enumerate}
    \item \textbf{Ziegler--Nichols (ZN)} \cite{ziegler1942optimum}: $K_p = 1.2\tau/(K\theta)$, $T_i = 2\theta$, $T_d = 0.5\theta$.
    \item \textbf{Cohen--Coon (CC)} \cite{cohen1953theoretical}: $K_p = \tau/(K\theta)\,(4/3 + r/4)$, $T_i = \theta(32 + 6r)/(13 + 8r)$, $T_d = 4\theta/(11 + 2r)$.
    \item \textbf{Tyreus--Luyben (TL)} \cite{tyreus1992tuning, luyben2007chemical}: $K_p = K_u/2.2$, $T_i = 2.2P_u$, $T_d = P_u/6.3$, with the ultimate gain and period of the FOPTD loop obtained exactly (Brent's method).
    \item \textbf{AMIGO} \cite{astrom2004revisiting, astrom2006pid}: $K_p = (0.2 + 0.45\,\tau/\theta)/K$, $T_i = \theta\,(0.4\theta + 0.8\tau)/(\theta + 0.1\tau)$, $T_d = 0.5\theta\tau/(0.3\theta + \tau)$.
    \item \textbf{SIMC} \cite{skogestad2003simple, grimholt2018optimal}: with $\tau_c = 0.40$~min, $K_p = \tau/(K(\tau_c + \theta))$, $T_i = \min(\tau, 4(\tau_c + \theta))$ and $T_d = 0.33\theta$.
    \item \textbf{Optimal ITAE}: the PID gains that minimize the ITAE of a $-2$~K step on the nonlinear plant and the benchmark valve subject to $M_s \le 1.6$ on the identified FOPTD loop \cite{garpinger2014performance}; the design and its cost are reported in Supplementary Table~\suppref{tab:constraint_cost}.
\end{enumerate}
The resulting gains are listed in Table~\ref{tab:testbed}. All six loops are reverse-acting ($K_p < 0$).
